Six out of seven G7 countries include standardized exams for university admission. 

This study documents the reordering of student's accessibility to elite university programs through a switch from standardized exam to non-standardized testing. I study this question by documenting grade inflation in the UK during the COVID-19 pandemic, caused by the cancellation of their standardized exam (A levels). Two months before the exam date, the British government banned in-person exams and allowed students to use their teachers predicted grade as their final A-level grade. Without any unified procedure for assessing grades, teachers provided grades using their personal knowledge and experience with the students. The resulting grade distribution moved dramatically upward, with the number of top grades increasing two times. The grade shift was mainly driven by grade inflation, as students had mostly finished their educational curriculum at the time of the cancellation, and students had already agreed with universities on the grade requirements for securing placements. 

I examine how the grading method change changed the probability of obtaining top grades across the students backgrounds. I focus on previous academic achievements, demographic attributes (race, gender, and parental income), subject content, and school of attendance and isolate the effects for each group. 

Firstly, I investigate whether schools assign grades to student according to the students previous academic standings. As British students typically take standardized national exams (GCSE) before entering higher school that strongly predicts achievement levels in the standardized exam, I examine whether the monotone relationship between grades and final grade held in the non-standardized exams. 

Secondly, I examine whether students' background demographics, such as gender and race, contributed to the grading. I focus on the role of race and gender on the students' rank within the entire distribution as well as within the students school and subject courses. I examine whether students parental backgrounds or their domicile country contributed to grading as well. 

Thirdly, I document whether grading patterns differed between different subject types. The UK allows university applicants to select 3 subjects among 100 subjects for the university admission test. Subjects ranges from quantitative subjects, such as mathematics and physics, to more qualitative subjects such as drama studies or arts. Teachers may hesitate to assign lower grades to students in non-quantitative subjects, as test answers include more discretion on the reviewers. Also subjects also differ in their importance on securing placements at universities, as elite universities in the UK prefer their applicants to take subjects from a pool of subjects (i.e. \emph{facilitating subjects}). I test whether higher increments in grade assignments were concentrated within competitive subjects. 

Finally, I examine whether schools differed in their inflating their pupils grades by their source of funding, selectivity of entrants, and discretion over choosing their curriculum. Approximately 7\% of university applicants prepare for their entry exams in privately funded schools that compete with each other schools by student placements. As private schools well-being is more closely attached to their pupil's placement outcomes, private school may have additional incentive to report higher grade than a student deserves with little respect to academic integrity. I also examine the extent of the increments across the schools with different value added. 

[Please Ignore] In the final part of the paper, I examine how the altered grade distribution changed the composition of students at universities. Despite the accusations that private schools inflated grades more than public schools, the school's impacts on compositions on students in university programs might be limited due to the students' already high chances of securing a placement in an elite program. On the other hand, students from public schools increase their representation in elite programs due to the large volume of students on the margins of receiving a high grade. Furthermore, the different extent of grade inflation across student groups and subject choice within state school may also reflect into the student composition at universities. I also take into account application patterns of different student groups, as previous literature documents underrepresented students apply much cautiously than other students.

\

I estimate grade inflation for each student using administrative data on entry levels A-level. The British higher education system is well suited for measuring grade inflation, as all students take standardized exams (GCSE) when they start their 16-18 education. The standardized measure of students academic performance allows me to compare the increments in their grade on a common measure of academic ability for the entire population of higher education applicants in the UK. Unsurprisingly, the GCSE scores accurately predict the results of the student grade A-levels that I test using an event-study test. 

I obtain the causal estimates of increments in the grade received by using an interrupted time series design that leverages the sudden announcement of the policy and the short time window for teachers to provide their predictions. Specifically, I measure grade inflation as the increments in the probability of the student receiving a top grade. Using administrative data on university application to the UK, I traces the grade outcomes of the universe of students and document the upward shift in final grade while conditioning the movements on the students' background or their previous academic performance. I include school-level and subject level control to document whether the resorting of students along the grade scale was drive by within-school or across schools. As the standard maximum-likelihood estimators are noisy with highe--dimensional models, I use ridge regression on a high-dimensional logit model to control for the school specific inflation levels.  The method recovers a rich set of estimates which I use to compare the degree of grade reassignment across a vast range of student backgrounds. 

[Please Ignore] In the final part of my paper, I develop a quantitative framework to compare the composition of student demographic group across universities. Following the general equilibrium framework proposed by Epple, Romano, and Sieg (2006), the model is calibrated to the exact demographic distribution of the student body in the UK, and capacity constraints of each British University. The model assigns grades differently based on the grading scheme; in which the non-standardized grading scheme assignment is based on the students school of attendance, their demographics, and subject. Student choose which university to apply to based on the rank of the universities. Universities set grade thresholds to optimize profits while keeping the volume of incoming cohorts under the capacity limit. Universities differ in their marginal costs for over/under accepting students. I additionally estimate the marginal costs by regressing the degrees of threshold adjustment in 2021 on the degree of over-acceptance in 2020 for each universities. 

\

I provide evidence of grade inflation after removing the standardized exams. Substantial heterogeneity arises in the attributes of the student, the school, and the exam subject. Conditional on the academic standing of the student, the odds increase is larger for students with higher academic standings. A student with the second highest academic standing (GCSE = 8) increased their odds of receiving an A/A* by 26 percent, while students with the lowest grade (GCSE = 3) decreased their odds by 50 percent. The increase/decrease in their odds is proportional to the students' achievement scores. In probability terms, a student with good academic standing with 50\% chances of receiving an A/A increased their probability to 58 \%, while a student with worse academic standing but with 50 \% on receiving A/A* reduced their chances of receiving A/A* to 35\%. The patterns in grade assignment and student achivements is robust to including subject level controls in the regression. This suggest that the ordering of students reflected the students previous academic achievements.

I find that female students increased their odds of receiving high grades. Female students increased their chances of receiving A / A * by 5 percentage points more than male students. The estimate is robust to inclusion of school-level controls, subject controls, and other demographic controls. This indicates that the gender inequality in grade achievement has decreased, as female students had a lower probability of obtaining high grades than male students on standardized exams. However, I do not find evidence on difference grades assigned by the ethnicity of the student, measured by white and nonwhites. Furthermore, income, measured by the occupation class of the parents, did not influence the final grades. I do not observe any patterns between income levels and the levels of grades assigned. However, the economic status of the permanent address of the students exhibits a positive relationship, as students from less deprived areas and students from higher college attainment population areas increases their chances in receiving higher grades. As these estimates are obtained from a school-subject levels regression, my results imply that students were graded differently based on their background demographics. 

I find that grade inflation was more concentrated in subjects with less quantitative content and less competitive subjects. Non-quantitative subjects were 1.5 percentage more likely to receive higher grades than quantitative subjects. In addition, less competitive subjects (non-facilitating subjects) were 0.5 percentage more likely to receive higher grades, even after controlling for the quantitative content within the non-facilitating subjects. 

I find evidence that independent schools substantially inflated grades more than publicly funded schools. Independent schools inflated grades by 7 percent more than publicly funded schools in 2021, while grade assignment patterns were similar between public and privately funded schools in 2020. I also find a negative relationship between school quality and the extent of inflation, as lower ranked schools inflate more to compensate for their lack of educational inputs. The negative relation ship is stronger for privately funded schools as the final probability of receiving top grade between the best and worst private schools significantly shrinks.

[Please ignore] My quantitative model predicts university composition of incoming students changes in particular student's originating schools and average academic performance in their high school exams. Taking the acceptance threshold as constant, the share of students from private schools \emph{decreases} by [] percent. In contrast, students in public universities increase by [] percent. This reflects the underlying connection between elite programs and private school institutions, private school students occupying almost [] percent of the student body in elite programs. In contrast, the share of students from public schools in elite programs increases by [] percent. The shift of students from public schools coincides with changes in demographic composition in elite universities, as the share of white and female students increases by [] percent in elite programs. In contrast, the share of white and female students in less elite programs is reduced by [] percent. 

My results indicates that the removal of standardized exams shifts grades unevenly upwards, but the magnitude of the increments in grades do not necessary improve access to elite universities. The nonstandardized grading scheme rescales the preexisting measure of academic performance of students, with the mean of the grade distribution moving upward. The resulting grade scheme is top-coded, failing to signal student quality to universities, especially for students at the top of the distribution. In addition, the upward shift in the grade distribution is accompanied by a reshuffling of achievement levels on their final exams across demographic groups of students. Student demographic groups sort to different schools (e.g. private or public) and their subject choice (e.g. quantitative or nonquantitative), which all receives different extent of grade inflation leading to a overall reshuffling of top grades. My results indicate that private schools unilaterally inflated their students' grades, indicating an increase in the number of students educated in private schools at the top of the grade scale. Female and white students moved upwards then other student groups, due to choosing exam subjects with more teacher leniency on grading. 

My paper points to how integrating non-standardizing testing into university admission changes the composition of universities with long-lasting impact on the graduates labor market outcomes.
\clearpage